% Compile with LuaLaTeX: `lualatex writeup_updown.tex`
\documentclass[11pt, a4paper]{article}
\usepackage{fontspec}
\usepackage[english]{babel}
\usepackage[a4paper,top=2cm,bottom=2cm,left=3cm,right=3cm,marginparwidth=1.75cm]{geometry}
\usepackage{amsmath, amssymb, amsthm}
\usepackage{graphicx}
\usepackage[colorlinks=true, allcolors=blue]{hyperref}

\def \N{{\mathbb Z}_{\geq 0}}
\def \Npos{{\mathbb Z}_{>0}}
\def \Z{\mathbb Z}
\linespread{1.2}
\setlength{\parindent}{0em}

\begin{document}

\section{Claude's review notes}

I went through the writeup once, cross-referenced every formula
against \texttt{crates/semigroup\_math} and the property assertions in
\texttt{tests/integration.rs::test\_up\_downs}, applied the fixes
directly in the following sections, and switched the preamble to a
LuaLaTeX-friendly setup. This section records what was changed and
what (if anything) is still open.

\subsection*{Math edits applied}

\begin{enumerate}
  \item \textbf{$r_i$ sign}. The original definition
        $r_i := (w_\mu - (w_{\mu-i}+w_i))/m$ contradicted the
        explanation on the same line ("$f = w_{\mu-i}+w_i-(r_i+1)m$",
        which solves to the opposite sign) and disagreed with
        \verb|Semigroup::r_i| in \texttt{semigroup.rs:192} (which is
        the Kunz coefficient $c(i, \mu-i)$). Now reads
        $r_i := (w_i + w_{\mu-i} - w_\mu)/m$.

  \item \textbf{$V(S)$}. Was $\{f-m+1, \dots, f-m\}$ (endpoints
        disagreed). Now $\{f-m+1, \dots, f-1\}$, matching
        \verb|v_in_s()| in \texttt{semigroup.rs:264}.

  \item \textbf{$\rho$}. Was $\min_i r_i$, which collapses to $0$
        because $r_\mu = 0$ for every numerical semigroup. Now
        $\rho(S) := \min_{i \neq \mu} r_i$, matching
        \verb|min_ri()| in \texttt{semigroup.rs:207}.

  \item \textbf{Apéry indexing}. "$w_i = i \bmod m$" tightened to
        "$w_i \equiv i \pmod m$", since $w_i$ is the unique Apéry
        element with residue $i$, not the residue itself.

  \item \textbf{Redundant clause in $\mathcal{D}_{total}$}. "$w_i =
        f+m$" together with $i = \mu$ is automatic; the
        load-bearing condition is $V(S) \subseteq S$. Now only the
        latter remains.

  \item \textbf{Subscript in \S~Correspondence}. The families are
        indexed $(i, m)$ throughout \S~UpDown, but
        \S~Correspondence had drifted to $(i, \mu)$. Restored to
        $(i, m)$.

  \item \textbf{Symmetric exclusion argument}. The original
        sentence justified excluding symmetric semigroups from
        $\mathcal{A}_{total}$ by appealing to $r = 0$, but that
        alone doesn't prevent $w_\mu$ from being an atom. Replaced
        with the actual reason: $\rho(S) = 0$ forces
        $w_i + w_{\mu-i} = w_\mu$ for some $i \neq \mu$, so $w_\mu$
        is decomposable and therefore not an atom; hence symmetric
        semigroups give an empty $\mathcal{A}_{total}$.

  \item \textbf{$a_\mu$ vs.\ $w_\mu$}. The body uses $w_\mu$
        throughout; one stray $a_\mu$ near the end of
        \S~Correspondence has been normalised.

  \item \textbf{Closing paren on the symmetric note}. "($r=0$ can
        never \dots)" $\to$ "($r=0$) can never \dots".

  \item \textbf{$S_+$}. Added a one-line declaration $S_+ := S
        \setminus \{0\}$ where the atom definition uses it.
\end{enumerate}

\subsection*{Typos / preamble cleanup}

Section heading "Corrospondence" $\to$ "Correspondence". "just just"
$\to$ "just". "relfected" $\to$ "reflected". "and$f$" got its space
back. "would us now require" $\to$ "would now require". "$\mod i$"
became "in residue class $i \pmod m$". The babel locale switched
from German to English (the prose is English). Duplicate
\verb|\usepackage| calls for \verb|amsmath|, \verb|amssymb|,
\verb|graphicx|, and \verb|geometry| collapsed to one each;
\verb|eurosym| was unused and is gone. \verb|inputenc| /
\verb|fontenc| removed (Lua\TeX{} reads UTF-8 natively); replaced
with \verb|fontspec|. The geometry options switched from
\verb|letterpaper| to \verb|a4paper| to match the document class.

\subsection*{Code block (\S~Code)}

Replaced the older verbatim with the current cleaned
\verb|test_up_downs| at
\texttt{crates/semigroup\_math/tests/integration.rs:531} (early-return
on $m=1$, formatted assertion messages, doc comment in the source).

\subsection*{What checks out cleanly}

The four $(\Delta g, \Delta \sigma, \Delta r)$ formulas all match the
runtime invariants:
\begin{itemize}
  \item $\mathcal{D}_{flip}$: \verb|down.r + 2 == s.r| and
        \verb|down.count_set == s.count_set + 1|.
  \item $\mathcal{D}_{total}$: \verb|down.r + 2 == s.r + s.m| and
        \verb|down.count_set + s.m == s.count_set + 1|.
  \item $\mathcal{A}_{flip}$: \verb|up.r == s.r + 2| and
        \verb|up.count_set + 1 == s.count_set|.
  \item $\mathcal{A}_{total}$: \verb|up.r + minr*(s.m-2) == s.r| and
        \verb|up.count_set == minr*(s.m-1) + s.count_set|.
\end{itemize}
The "trivial" identities $r + \sigma = g$, $f + 1 + r = 2g$, and $t -
1 \leq r$ are asserted on every fixture (\texttt{integration.rs:358-385}).

\subsection*{Notation $\leftrightarrow$ codebase}

\begin{center}
\begin{tabular}{l|l}
LaTeX & Code \\\hline
$m,\,f,\,\mu$           & \verb|s.m|, \verb|s.f|, \verb|s.mu| \\
$g,\,\sigma,\,r$        & \verb|s.count_gap|, \verb|s.count_set|, \verb|s.r| \\
$r_i,\,\rho$            & \verb|s.r_i(i)|, \verb|s.min_ri()| \\
$w_i,\,\operatorname{Ap}(S)$ & \verb|s.apery_set[i]|, \verb|s.apery_set| \\
$\operatorname{RG}(S)$  & \verb|s.blob()| \\
$V(S) \subseteq S$      & \verb|s.v_in_s()| \\
atom                    & element of \verb|s.gen_set| \\
$S^\downarrow$ (descent) & \verb|s.toggle(w - s.m)| \\
$S^\uparrow$ (ascent)    & \verb|s.toggle(w)| \\
$t$                     & \verb|s.t| (Rust) / \verb|s.type_t| (JS)
\end{tabular}
\end{center}

Not used in this note but tracked in the app: $r_a$
(\verb|s.ra|, reflected Apéry), $e_s$ (\verb|s.es|, small atoms $g <
f-m$), $r_l$ (\verb|s.rl|, large reflected gaps $L \in (f-m, f)$).

\subsection*{Still open (your call)}

\begin{itemize}
  \item \S~Correspondence is doing two things at once: (a) when is
        $D \circ A = \mathrm{id}$ on $\mathcal{A}_{total}$, (b) when
        is $A \circ D = \mathrm{id}$ on a restricted
        $\mathcal{D}_{total}$. Splitting these explicitly would
        help; I left the text as-is to preserve your voice.
  \item The closing question "we need to restrict to semigroups in
        $\mathcal{D}_{total}(i,m)$ where $w_\mu$ is not an atom?"
        ends with a question mark. Fine for a chat-with-friends
        register, but if you intend it as a conclusion the wording
        could be sharpened.
\end{itemize}

\section{Setup}

Let $S$ be a numerical semigroup, with $f,\,m,\,g,\,\sigma$ as usual.
Write $S_+ := S \setminus \{0\}$.
\begin{itemize}
    \item $\operatorname{Ap}(S):=\{s\in S \mid s-m \notin S\}$,
    $\operatorname{RG}(S):=\{L\notin S \mid f-L \notin S\}$.
    \item $a$ is an \textit{atom} iff $a\in S\setminus (S_+ + S_+)$
    (the common name "minimal generator" actually means an element of
    the unique minimal generating set).
    \item $\operatorname{Ap}(S):=\{w_0,w_1,\dots,w_{m-1}\}$ where
    $w_i \equiv i \pmod m$.
    \item $\mu \in \{1,\dots,m-1\}$ is the number $\mu = f \bmod m$,
    equivalently $f+m = w_\mu$.
    \item For every $1\leq i \leq m-1$ let
    \[
        r_i := \frac{w_i + w_{\mu-i} - w_\mu}{m}
    \]
    (the number of reflected gaps in residue class $i \pmod m$). Just
    look at $w_\mu - m = f = w_{\mu-i} + w_i - (r_i+1)m$, i.e.\
    distribute the $r_i+1$ gaps between $w_\mu$ and $w_i + w_{\mu-i}$
    along the residue class $\mu$.
    \item $V(S):=\{f-m+1,\,f-m+2,\dots,f-1\}$, the
    \textit{interval before $f$}.
    \item $\rho(S):=\min_{i\neq\mu} r_i$ and
    $r = \sum_i r_i = \#\operatorname{RG}(S)$.
    \item Trivial: $r+\sigma=g$, $f+1+r=2g$, $t-1\leq r$.
\end{itemize}

\section{UpDown}

Now define four families of semigroups, for a fixed $m>2$ and
an $i \in \{1,\dots,m-1\}$, indicating by $'$ the new count in
$S^\uparrow$ resp.\ $S^\downarrow$:
\begin{enumerate}
    \item $\mathcal{D}_{flip}(i,m)$: $i\neq \mu$, $S$ and
    $f<w_i<f+m$.
    \[S^\downarrow:=S\cup\{w_i-m\}\quad g'=g-1\quad \sigma'=\sigma+1\quad r'=r-2\]
    \item $\mathcal{D}_{total}(i,m)$: $i=\mu$ and
    $V(S)\subseteq S$ (then $w_\mu = f+m$ automatically).
    \[S^\downarrow:=S\cup\{w_\mu-m\}\quad g'=g-1\quad \sigma'=\sigma-(m-1)\quad r'=r+(m-2)\]
    \item $\mathcal{A}_{flip}(i,m)$: $i\neq \mu$, $S$ and $w_i$ is an
    atom, $w_i \in V(S)$.
    \[S^\uparrow:=S\setminus\{w_i\}\quad g'=g+1\quad \sigma'=\sigma-1\quad r'=r+2\]
    \item $\mathcal{A}_{total}(i,m)$: $i=\mu$, $S$ and $w_\mu$ is an
    atom; set $\rho=\rho(S)$.
    \[S^\uparrow:=S\setminus\{w_\mu\}\quad g'=g+\rho \quad \sigma'=\sigma+\rho(m-1)\quad r'=r-\rho(m-2)\]
\end{enumerate}
All these claims can be proven by carefully checking the change in
the semigroup --- and the property tests in \texttt{tests/integration.rs}
verify each $(\Delta g, \Delta \sigma, \Delta r)$ on every fixture in
\texttt{tests/data/check\_fixtures.csv}.

\section{Correspondence}

If we want a bijection between
\[
\mathcal{D}_{total}(i,m) \longleftrightarrow \mathcal{A}_{total}(i,m)
\]
we obviously need $\rho(S)=1$ for the semigroup in
$\mathcal{A}_{total}(i,m)$; this is a necessary condition. The reason
is simple: if $\rho > 1$, then $w_\mu+m$ cannot be written as a sum of
Apéry elements, so it will not be in $S^\uparrow$. Note that symmetric
semigroups ($r=0$) can never be in $\mathcal{A}_{total}(i,m)$: $r=0$
forces $\rho=0$, which forces some $w_i + w_{\mu-i} = w_\mu$ with
$i \neq \mu$, so $w_\mu$ is decomposable and therefore not an atom.
\vskip5mm
Viewed from below we need
\[
g + \rho -1 = g
\qquad
\sigma + \rho (m-1) - (m-1) = \sigma
\qquad r-\rho(m-2)+(m-2)=r
\]
It is tempting to define adding $f$ or removing $f+m$ by maps
\[
D: \mathcal{D}_{total}(i,m) \rightarrow \mathcal{A}_{total}(i,m)
\qquad D(S)=S\cup\{f(S)\} = S^\downarrow
\] and \[
A: \mathcal{A}_{total}(i,m) \rightarrow \mathcal{D}_{total}(i,m)
\qquad A(S) = S \setminus \{w_\mu(S)\} = S^\uparrow
\]
This would be great --- at least for the $r_i$, because in this case
we would have $r_i(D(S))=r_i(S)+1$ and $r_i(A(S))=r_i(S)-1$, which
would now require that the semigroup in $\mathcal{D}_{total}(i,m)$ has
$\rho=0$. That means $w_\mu(S^\uparrow)$ is decomposable and not an
atom. It is obvious that $V(A(S))\subseteq A(S)$ --- so in order to
get a well-defined bijection, we need to restrict to semigroups in
$\mathcal{D}_{total}(i,m)$ where $w_\mu$ is not an atom?

\section{Code}
\begin{verbatim}
fn test_up_downs(s: &Semigroup) {
    if s.m == 1 {
        return;
    }
    for i in 1..s.m {
        let w = s.apery_set[i];

        // Descent: w-m is a gap; adding it lifts the membership of w-m.
        if w > s.f && s.f > s.m {
            assert!(!s.element(w - s.m));
            assert!(w <= s.f + s.m);
            assert!(
                w == s.f + s.m || s.is_reflected_gap(w - s.m),
                "reflected {w} at {i}"
            );
            let down = s.toggle(w - s.m);
            assert_eq!(down.count_gap + 1, s.count_gap);
            if w == s.f + s.m && s.v_in_s() {
                assert_eq!(down.r + 2, s.r + s.m, "r change w==f+m for i={i}");
                assert_eq!(down.count_set + s.m, s.count_set + 1, "count_set");
            }
            if w < s.f + s.m {
                assert_eq!(down.r + 2, s.r, "r change w<f+m for i={i}");
                assert_eq!(down.count_set, s.count_set + 1, "count_set");
            }
        }

        // Ascent: w must already be a minimal generator.
        if s.gen_set.contains(&w) {
            let up = s.toggle(w);
            if w < s.f && w + s.m > s.f {
                assert_eq!(
                    up.count_gap,
                    s.count_gap + 1,
                    "removing atom generates gap for {i} and {w}"
                );
                assert!(up.apery_set.contains(&(w + s.m)));
                assert_eq!(up.apery_set[i], w + s.m);
                assert_eq!(up.r, s.r + 2, "r change w<f for i={i}");
                assert_eq!(up.count_set + 1, s.count_set);
            }
            if w == s.f + s.m {
                let minr = s.min_ri();
                assert_eq!(
                    up.count_gap,
                    s.count_gap + minr,
                    "removing atom generates minr gap for {i} and {w}"
                );
                assert!(up.apery_set.contains(&(w + minr * s.m)));
                assert_eq!(up.apery_set[i], w + minr * s.m);
                assert_eq!(up.r + minr * (s.m - 2), s.r, "r change w==f+m for i={i}");
                assert_eq!(
                    up.count_set,
                    minr * (s.m - 1) + s.count_set,
                    "sigma change w==f+m for i={i}"
                );
            }
        }
    }
}
\end{verbatim}

\end{document}
